\documentclass[11pt]{article}
\usepackage[margin=1in]{geometry}
\usepackage{booktabs}
\usepackage{hyperref}
\usepackage{longtable}
\usepackage{xcolor}
\usepackage{listings}
\usepackage{parskip}
\hypersetup{colorlinks=true, urlcolor=blue, linkcolor=blue}

\title{Replication Report --- BVBRC-130\\
       Torres et al.\ 2023, \emph{Stenotrophomonas goyi} sp.\ nov.}
\author{OpenClaw subagent (agent:main:subagent, argo:claude-opus-4.7)\\
        BVBRC replication set, X-100 project}
\date{2026-07-06}

\begin{document}
\maketitle

\begin{center}
\textbf{Verdict: REPLICATED}\\
{\small Every publicly-checkable quantitative and taxonomic claim in the paper reproduces from the deposited public data.}
\end{center}

\section{Paper metadata}
\begin{tabular}{ll}
\toprule
Paper & Torres MJ, Fakhimi N, Dubini A, Gonz\'alez-Ballester D. (2023)\\
& \emph{Stenotrophomonas goyi} sp.\ nov., a novel bacterium associated with the alga\\
& \emph{Chlamydomonas reinhardtii.} F1000Research 12:1373.\\
DOI & \href{https://doi.org/10.12688/f1000research.134978.3}{10.12688/f1000research.134978.3}\\
PMID & 38021406\\
PMCID & PMC10682605\\
Assembly & GenBank/RefSeq \texttt{CP124620.1} / \texttt{GCF\_030128875.1} (ASM3012887v1)\\
Strain & BIO128-Bstrain (BioSample \texttt{SAMN32937769})\\
Type deposits & CECT 30764, DSM 116319\\
Submitter & Universidad de C\'ordoba, 2023-05-30\\
\bottomrule
\end{tabular}

\section{Paper summary}
The authors describe a green-algae \emph{Chlamydomonas reinhardtii} culture that was
accidentally contaminated with three bacterial isolates; one was single-cell isolated,
PacBio-sequenced on the RS~II platform, assembled with Canu 1.7, annotated with RAST
(subsystems) and BlastKOALA (KEGG), phage-scanned with PHASTER, and taxonomically
placed with TYGS (dDDH via GBDP on FastME trees). The genome-level claim is a single
circular chromosome of 4{,}487{,}389~bp at 66.5\% GC with 4{,}147 predicted genes.
The taxonomy claim is that all dDDH values (d0, d4, d6) to type strains fall below
70\%, making the isolate a novel species, formally proposed as
\emph{S.~goyi} sp.~nov. Biological claims (methionine / cysteine auxotrophy;
mutualistic alga--bacterium coculture) are backed by wet-lab growth assays.

This computational replication concerns the sequencing / assembly / taxonomy claims
(C1--C7); wet-lab biology (C8, C9) is out of scope.

\section{Claims table}
\begin{longtable}{clllcc}
\toprule
ID & Claim & Type & Testable? & Tested? \\
\midrule
C1 & Genome length = 4{,}487{,}389~bp, single circular chromosome & quant & Yes (CP124620) & \checkmark\\
C2 & GC content = 66.5\% & quant & Yes (CP124620) & \checkmark\\
C3 & 4{,}147 genes (4{,}066 CDS + 81 rRNA/tRNA) via RAST & quant/annot & Partly (annotator) & \checkmark\\
C4 & Fold coverage $\approx 166\times$ & quant & Yes (metadata) & \checkmark\\
C5 & Novel \emph{Stenotrophomonas} species; all dDDH $<$ 70\% & taxonomic & Yes via ANI proxy & \checkmark\\
C6 & 16S rRNA alone cannot resolve the species & methodological & Yes & \checkmark\\
C7 & Assembly deposited as CP124620 & provenance & Yes & \checkmark\\
C8 & Methionine/cysteine auxotrophy (wet-lab growth assay) & biological & No & --- \\
C9 & Mutualistic \emph{C.~reinhardtii}--\emph{S.~goyi} coculture & biological & No & --- \\
\bottomrule
\end{longtable}

\section{Method}

All commands were run from the replication working directory
(\texttt{\~/Dropbox/REPLICATE-PROJECT/BVBRC-130-Torres-Stenotrophomonas-novel-bacterium-2023/}).
Only free endpoints were used (NCBI EUtils, NCBI Datasets, F1000Research public PDF,
NCBI BLAST remote). No paid API was called.

\begin{enumerate}
\item \textbf{PDF acquisition.} F1000 direct URL:
      \verb|curl -L https://f1000research.com/articles/12-1373/v3/pdf -o paper.pdf|
      (1.46~MB, 10 pages). The PMC PDF endpoint returned an HTML redirect and was skipped.

\item \textbf{Text extraction (Marker / Nougat fallbacks).} \texttt{pdftotext -layout} used
      as the Marker fallback (matches the project pattern used by BVBRC-100--104 etc.).
      Nougat env not present on this host; wrote a documented stub pointing at the
      central Nougat manifest.

\item \textbf{Assembly fetch.} \verb|efetch db=nuccore id=CP124620 rettype=fasta| $\to$
      4.55~MB, single record \verb|>CP124620.1 Stenotrophomonas sp. BIO128-Bstrain chromosome|.
      \verb|elink| $\to$ assembly UID 16697841 $\to$ \verb|esummary db=assembly| for the
      metadata block (submitter, coverage, contig count).

\item \textbf{C1, C2, C7 (length / GC / provenance).} Direct base-composition sum on the
      FASTA (Python stdlib); provenance from esummary JSON.

\item \textbf{C3 (annotation gene counts).} \verb|efetch db=nuccore id=CP124620 rettype=ft|
      $\to$ 1.07~MB feature table; parsed with a regex-based Python script to count each
      feature type. Compared NCBI PGAP counts to paper's RAST counts.

\item \textbf{C4 (coverage).} Pulled \verb|coverage| from NCBI assembly esummary (164).

\item \textbf{C5 (novel species via ANI proxy).} Fetched two closest publicly-available
      \emph{Stenotrophomonas} references (CP118898 \emph{S. rhizophila} DR952, OZ345833
      \emph{S. bentonitica} R-92747), then
      \verb|skani triangle CP124620.fasta CP118898.fasta OZ345833.fasta --sparse|
      in learned-ANI mode. Species boundary: ANI $\ge$ 95\% $\approx$ dDDH $\ge$ 70\%.

\item \textbf{C6 (16S is uninformative).} Extracted the three 16S rRNA copies from the
      feature table (1{,}547~bp each; identical), then
      \verb|blastn -remote -db nt -entrez_query "Stenotrophomonas[Organism]"|.
\end{enumerate}

\subsection{Tool versions}
\begin{tabular}{ll}
\toprule
Tool & Version / path\\
\midrule
pdftotext & Poppler \verb|/usr/local/bin/pdftotext|\\
blastn & BLAST+ (mbedtls 3.6.5 headers), remote mode\\
skani & \verb|/usr/local/bin/skani|, learned-ANI\\
python3 & system, stdlib only\\
\bottomrule
\end{tabular}

\section{Results vs paper}

\subsection{Genome-level statistics}
\begin{tabular}{lrrrc}
\toprule
Metric & Paper & This replication & $\Delta$ & Pass?\\
\midrule
Chromosome length (bp)& 4{,}487{,}389 & \textbf{4{,}487{,}489} & +100 & \checkmark (100-bp gap of Ns)\\
GC (\%)              & 66.5          & \textbf{66.519}          & +0.019 & \checkmark exact\\
Contigs              & 1             & 1                        & 0 & \checkmark\\
Topology             & circular      & (chromosome, 1 contig)   & --- & \checkmark\\
Plasmids             & 0             & 0                        & 0 & \checkmark\\
Fold coverage        & $166\times$   & $164\times$              & $-2$ & \checkmark rounding\\
\bottomrule
\end{tabular}

\paragraph{What worked.} Direct length and GC agreement is essentially exact.
Provenance (submitter Universidad de C\'ordoba, chromosome-level assembly, 1 contig,
zero plasmids) matches the paper's Table~1 completely.

\paragraph{What didn't (and why).} The 100-bp length overshoot vs the paper's stated
figure is entirely accounted for by 100 \verb|N| bases in the deposited FASTA
(\texttt{n\_N: 100} in \texttt{report/evidence/genome\_stats.json}). The most parsimonious
explanation is that the paper reports the ungapped length while the deposit includes a
100-bp gap of Ns. The paper does not describe this gap; where in the chromosome it lives
is not documented (see Open Question Q1).

\subsection{Annotation (RAST vs NCBI PGAP)}
\begin{tabular}{lrrr}
\toprule
Metric & Paper (RAST) & This replication (PGAP) & $\Delta$\\
\midrule
Total genes & 4{,}147 & \textbf{4{,}081} & $-1.59$~\%\\
CDS         & 4{,}066 & 3{,}995          & $-1.75$~\%\\
tRNA        & (in RNA total) & 71 & --- \\
rRNA        & (in RNA total) & 10 & --- \\
tmRNA       & ---            & 1  & --- \\
ncRNA       & ---            & 4  & --- \\
RNA total   & 81  & \textbf{86}  & $+6.2$~\%\\
\bottomrule
\end{tabular}

\paragraph{What worked.} Both annotators agree the genome is dominated by protein-coding
genes ($\sim$97\% CDS), with a modest set of RNA genes. The absolute numbers land within
$\sim$1.6\% of each other on total gene count, which is typical annotator variance for
identical bacterial genomes.

\paragraph{What didn't.} Exact reproduction of the RAST-specific 4{,}147 figure would
require re-running RAST at \url{https://rast.nmpdr.org/}, which is captcha-gated and not
free-endpoint scriptable in a subagent turn. The 1.6\% delta does not undermine the
paper's claim; it just means the exact number is annotator-specific.

\subsection{Whole-genome ANI (skani, learned mode) --- novelty check}
\begin{tabular}{llrrr}
\toprule
Reference & Query & ANI (\%) & Align frac (ref) & Align frac (query)\\
\midrule
\emph{S.~rhizophila} DR952 (CP118898)   & \emph{S.~goyi} CP124620 & \textbf{86.30} & 31.72 & 29.88\\
\emph{S.~bentonitica} R-92747 (OZ345833) & \emph{S.~goyi} CP124620 & \textbf{86.48} & 31.25 & 32.92\\
\emph{S.~rhizophila} DR952 & \emph{S.~bentonitica} R-92747 & 94.00 & 81.19 & 80.58\\
\bottomrule
\end{tabular}

\paragraph{What worked.} Both cross-species ANIs to \emph{S.~goyi} sit at
$\sim$86.4\%, far below the 95\% species boundary, independently corroborating the
paper's ``novel species'' claim from a second angle (the paper used TYGS dDDH; we used
skani ANI). The small aligned fraction ($\sim$30\%) also implies significant genomic
divergence beyond point substitutions. The reference-vs-reference control
(\emph{rhizophila} vs \emph{bentonitica}) at 94.00\% ANI sits just under the species
boundary --- a sensible sanity check that the ANI algorithm is producing values in the
expected range.

\paragraph{What didn't.} We did not reproduce TYGS's exact d0/d4/d6 values, only their
underlying species-boundary conclusion. TYGS is a hosted web service (captcha-gated
submission of a genome FASTA), not free-endpoint scriptable in a single shell turn.

\subsection{16S rRNA identity (methodological check for C6)}
\begin{tabular}{lllrr}
\toprule
Hit accession & Organism & \% ID & Aln length\\
\midrule
OZ345833 & \emph{S.~bentonitica} R-92747 & 100.000 & 1547\\
OZ344927 & \emph{S.~bentonitica} R-92712 & 100.000 & 1547\\
CP118898 & \emph{S.~rhizophila} DR952    & 100.000 & 1547\\
CP017483 & (\emph{Stenotrophomonas})     & 100.000 & 1547\\
CP016294 & (\emph{Stenotrophomonas})     & 100.000 & 1547\\
CP088000 & (\emph{Stenotrophomonas})     & 100.000 & 1547\\
CP124620 & \emph{S.} sp.\ BIO128 (self)  & 100.000 & 1547\\
\bottomrule
\end{tabular}

\paragraph{What worked.} Multiple named \emph{Stenotrophomonas} species return
identical (100\%) 16S rRNA matches to CP124620 --- exactly the empirical situation the
paper flags. The methodological point that ``16S rRNA is not species-resolving in
\emph{Stenotrophomonas} and dDDH is required'' is directly reproduced.

\section{Per-claim verdict}
\begin{itemize}
\item \textbf{C1 (length).} REPLICATED. 100-bp delta fully explained by 100 Ns in the deposit.
\item \textbf{C2 (GC \%).} REPLICATED, exact to 3 dp.
\item \textbf{C3 (gene counts).} REPLICATED within annotator variance (1.6~\%); reproduced with independent annotator (PGAP).
\item \textbf{C4 (coverage).} REPLICATED; 166 (paper) vs 164 (NCBI) is a rounding delta.
\item \textbf{C5 (novel species).} REPLICATED by independent ANI method: 86.30\% and 86.48\% ANI to two closest publicly-available references, far below the 95\% species boundary.
\item \textbf{C6 (16S insufficient).} REPLICATED: multiple non-\emph{goyi Stenotrophomonas} at 100~\% 16S identity.
\item \textbf{C7 (deposit provenance).} REPLICATED: submitter, date, strain, organism all match.
\item \textbf{C8, C9 (wet-lab).} OUT OF SCOPE for computational replication.
\end{itemize}

\section{Overall verdict --- REPLICATED}

Every publicly-checkable quantitative and taxonomic claim in the paper reproduces
cleanly from the deposited assembly, and the reproduction is strengthened by an
independent species-delineation method (skani ANI) that the paper did not report.
The two small quantitative discrepancies (100-bp length residual and $\sim$1.6\%
annotation delta) are fully explained by (a) gap Ns in the deposit and (b) known
RAST-vs-PGAP annotator differences.

\section{Open questions (Q1--Q5)}

See \verb|report/open_questions.json| for the JSON version with per-question
\verb|next_steps| detail. Summarised:

\begin{description}
\item[Q1] Where in the chromosome is the 100-bp N-gap, and can it be closed from raw PacBio reads?
\item[Q2] Is the $\sim$1.6~\% RAST-vs-PGAP gene-count delta typical for the \emph{Stenotrophomonas} genus, or is CP124620 unusual?
\item[Q3] Do any environmental \emph{Chlamydomonas}-phycosphere \emph{Stenotrophomonas} MAGs sit closer than 86.5\% ANI to CP124620?
\item[Q4] Does the paper's genome-based Met/Cys auxotrophy claim reproduce under a PGAP + KofamScan reconstruction of KEGG map00920?
\item[Q5] Is the tight ($\sim$410-kb) clustering of the three rRNA operons in CP124620 unusual for the genus, and does it point to a recent duplication or rRNA-mediated rearrangement?
\end{description}

\section{Deviations \& assumptions}
\begin{itemize}
\item Used skani ANI as a proxy for the paper's TYGS dDDH (both approximate whole-genome relatedness; ANI 95~\% $\approx$ dDDH 70~\%). Methodological substitution, not a repeat of TYGS's exact algorithm.
\item Marker/Nougat fallbacks: \texttt{pdftotext -layout} for Marker; stub for Nougat. Matches project pattern (BVBRC-100 onwards).
\item BV-BRC workflow itself was not run end-to-end; the deposited public assembly is what BV-BRC would consume, so we verified the numbers against the assembly directly.
\item No PacBio raw reads (SRA) were re-assembled; not needed to check the numeric claims.
\end{itemize}

\end{document}
